Кыргызстанда дизайн-студиянын кардарлары долбоорду көбүнчө 360\textdegree{} панорамалар жана рендерлер аркылуу бекитет. Камеранын бурчун дизайнер тандайт, ал эми кең бурчтуу камера жана <<идеалдуу>> жарык бөлмөлөрдү чындыгындагыдан чоңураак жана жарыгыраак кылып көрсөтөт. Натыйжада кардар бөлмөнүн өлчөмүн, жарыкты, эшик менен эмеректердин кантип ачыларын туура эмес баалайт, ал эми мындай каталар оңдоо иштеринин учурунда гана билинет.

Бул иште JK3DStudio VR сунушталат "--- 3ds Max программасында даярдалган дизайн-долбоорду интерактивдүү виртуалдык чындык (VR) презентациясына айландыруучу курал. Кардар келечектеги батирди 1:1 масштабда кыдырып, эшиктерди колу менен ачат, ар бир бөлмөдө жарыкты күйгүзөт, күндүн убактысын жана жарыктын түс температурасын өзгөртөт. Орус жана англис тилдеринде сүйлөгөн ЖИ-гид дизайнер бекиткен долбоордун паспортундагы фактыларды гана айтат.

Система үч бөлүктөн турат: аталыштардын стандарты боюнча сахнаны жарым-автоматтык түрдө даярдаган Unreal Engine 5.8 плагини, PICO 4 Ultra көз айнегине арналган VR-тиркеме жана сүйлөө текстин түзүү, фактыларды текшерүү жана үн берүү үчүн AI Gateway сервердик бөлүгү (FastAPI, PostgreSQL). Чет өлкөлөрдө мындай интерактивдүү турлар кол менен жасалат жана баасы 10\,000 АКШ долларынан башталат. Кыргызстанда дизайн-долбоорлор үчүн мындай кызмат табылган жок. Жаңылыгы "--- тез даярдоо, чыныгы масштаб жана жарык, ошондой эле жергиликтүү рынокко ыңгайлаштырылган, фактыларга таянган ЖИ-гид. Биринчи алдын ала коргоодо VR-тиркеменин иштеген прототиби жана текшерилген сервердик бөлүгү көрсөтүлөт.

\medskip\noindent\textbf{Негизги сөздөр:} виртуалдык чындык, интерьер дизайны, Unreal Engine, ЖИ-гид, чоң тил моделдери, Кыргызстан.
